
\documentclass{article}
\usepackage{colortbl}
\usepackage{makecell}
\usepackage{multirow}
\usepackage{supertabular}

\begin{document}

\newcounter{utterance}

\centering \large Interaction Transcript for game `referencegame', experiment `letter\_grids', episode 2 with Qwen3.6{-}35B{-}A3B{-}FP8{-}without{-}reasoning.
\vspace{24pt}

{ \footnotesize  \setcounter{utterance}{1}
\setlength{\tabcolsep}{0pt}
\begin{supertabular}{c@{$\;$}|p{.15\linewidth}@{}p{.15\linewidth}p{.15\linewidth}p{.15\linewidth}p{.15\linewidth}p{.15\linewidth}}
    \# & $\;$A & \multicolumn{4}{c}{Game Master} & $\;\:$B\\
    \hline

    \theutterance \stepcounter{utterance}  
    & & \multicolumn{4}{p{0.6\linewidth}}{
        \cellcolor[rgb]{0.9,0.9,0.9}{
            \makecell[{{p{\linewidth}}}]{
                \texttt{\tiny{[P1$\langle$GM]}}
                \texttt{Vou te mostrar três tabuleiros quadriculados, cada um de tamanho 5 por 5.} \\
\texttt{Os tabuleiros têm quadros vazios marcados por "▢" e quadros preenchidos marcados por "X".} \\
\texttt{Sua tarefa é gerar uma expressão de referência que melhor descreva o tabuleiro{-}alvo, de forma a distingui{-}lo dos outros dois tabuleiros distratores.} \\
\texttt{O primeiro tabuleiro é o tabuleiro{-}alvo, e os dois seguintes são distratores.} \\
\\ 
\texttt{Tabuleiro{-}alvo:} \\
\\ 
\texttt{X X X X X} \\
\texttt{X □ X □ X} \\
\texttt{X X X X X} \\
\texttt{X □ X □ X} \\
\texttt{X X X X X} \\
\\ 
\texttt{Tabuleiro distrator 1:} \\
\\ 
\texttt{X X X X X} \\
\texttt{X □ □ □ □} \\
\texttt{X X X X X} \\
\texttt{X □ □ □ □} \\
\texttt{X X X X X} \\
\\ 
\texttt{Tabuleiro distrator 2:} \\
\\ 
\texttt{X X X X X} \\
\texttt{X □ □ □ □} \\
\texttt{X □ □ □ □} \\
\texttt{X □ □ □ □} \\
\texttt{X X X X X} \\
\\ 
\texttt{Instrução: Descreva o tabuleiro{-}alvo.} \\
\texttt{Gere a expressão de referência começando com a etiqueta "Expressão: " para o tabuleiro{-}alvo indicado. Omita qualquer outro texto.} \\
            }
        }
    }
    & \\ \\

    \theutterance \stepcounter{utterance}  
    & \multicolumn{4}{p{0.6\linewidth}}{
        \cellcolor[rgb]{0.9,0.9,0.9}{
            \makecell[{{p{\linewidth}}}]{
                \texttt{\tiny{[P1$\rangle$GM]}}
                \texttt{Expressão: O tabuleiro{-}alvo possui uma estrutura de grade onde as linhas ímpares (1, 3, 5) são completamente preenchidas com "X", enquanto nas linhas pares (2, 4), apenas as colunas ímpares (1, 3, 5) contêm "X", e as colunas pares (2, 4) estão vazias ("▢").} \\
            }
        }
    }
    & & \\ \\

    \theutterance \stepcounter{utterance}  
    & & & \multicolumn{2}{p{0.3\linewidth}}{
        \cellcolor[rgb]{0.9,0.9,0.9}{
            \makecell[{{p{\linewidth}}}]{
                \texttt{\tiny{[GM$|$GM]}}
                \texttt{Expressão: O tabuleiro{-}alvo possui uma estrutura de grade onde as linhas ímpares (1, 3, 5) são completamente preenchidas com "X", enquanto nas linhas pares (2, 4), apenas as colunas ímpares (1, 3, 5) contêm "X", e as colunas pares (2, 4) estão vazias ("▢").} \\
            }
        }
    }
    & & \\ \\

    \theutterance \stepcounter{utterance}  
    & & \multicolumn{4}{p{0.6\linewidth}}{
        \cellcolor[rgb]{0.9,0.9,0.9}{
            \makecell[{{p{\linewidth}}}]{
                \texttt{\tiny{[GM$\rangle$P2]}}
                \texttt{Vou te mostrar três tabuleiros quadriculados, cada um de tamanho 5 por 5.} \\
\texttt{Os tabuleiros têm quadros vazios marcados por "▢" e quadros preenchidos marcados por "X".} \\
\texttt{Vou te dar também um expressão de referência que descreve um desses tabuleiros.} \\
\texttt{Sua tarefa é selecionar o tabuleiro que combina com essa expressão de referência.} \\
\texttt{Gere apenas o número (em forma de texto) do tabuleiro que combina com essa expressão de referência, selecionando primeiro, segundo ou terceiro.} \\
\\ 
\texttt{primeiro:} \\
\\ 
\texttt{X X X X X} \\
\texttt{X □ □ □ □} \\
\texttt{X □ □ □ □} \\
\texttt{X □ □ □ □} \\
\texttt{X X X X X} \\
\\ 
\texttt{segundo:} \\
\\ 
\texttt{X X X X X} \\
\texttt{X □ □ □ □} \\
\texttt{X X X X X} \\
\texttt{X □ □ □ □} \\
\texttt{X X X X X} \\
\\ 
\texttt{terceiro:} \\
\\ 
\texttt{X X X X X} \\
\texttt{X □ X □ X} \\
\texttt{X X X X X} \\
\texttt{X □ X □ X} \\
\texttt{X X X X X} \\
\\ 
\texttt{Expressão: O tabuleiro{-}alvo possui uma estrutura de grade onde as linhas ímpares (1, 3, 5) são completamente preenchidas com "X", enquanto nas linhas pares (2, 4), apenas as colunas ímpares (1, 3, 5) contêm "X", e as colunas pares (2, 4) estão vazias ("▢").} \\
\texttt{Pergunta: A qual dos tabuleiros essa expressão de referência se refere?} \\
\texttt{Comece com a etiqueta "Resposta: ", seguida da sua escolha. Omita qualquer outro texto.} \\
            }
        }
    }
    & \\ \\

    \theutterance \stepcounter{utterance}  
    & & & \multicolumn{4}{p{0.6\linewidth}}{
        \cellcolor[rgb]{0.9,0.9,0.9}{
            \makecell[{{p{\linewidth}}}]{
                \texttt{\tiny{[GM$\langle$P2]}}
                \texttt{Resposta: terceiro} \\
            }
        }
    }
     \\ \\

    \theutterance \stepcounter{utterance}  
    & & & \multicolumn{2}{p{0.3\linewidth}}{
        \cellcolor[rgb]{0.9,0.9,0.9}{
            \makecell[{{p{\linewidth}}}]{
                \texttt{\tiny{[GM$|$GM]}}
                \texttt{terceiro} \\
            }
        }
    }
    & & \\ \\

\end{supertabular}
}

\end{document}
